\documentclass[12pt, a4paper]{report}
\usepackage[utf8]{inputenc}
\usepackage{geometry}
\geometry{margin=1in}
\usepackage{hyperref}
\usepackage{titlesec}
\usepackage{enumitem}
\usepackage{xcolor}
\usepackage{setspace}
\usepackage{listings}
\usepackage{tcolorbox}
\setstretch{1.2}

\definecolor{codegray}{rgb}{0.5,0.5,0.5}
\definecolor{codepurple}{rgb}{0.58,0,0.82}
\definecolor{backcolour}{rgb}{0.95,0.95,0.92}

\lstdefinestyle{mystyle}{
    backgroundcolor=\color{backcolour},   
    commentstyle=\color{codegray},
    keywordstyle=\color{magenta},
    numberstyle=\tiny\color{codegray},
    stringstyle=\color{codepurple},
    basicstyle=\ttfamily\footnotesize,
    breakatwhitespace=false,         
    breaklines=true,                 
    captionpos=b,                    
    keepspaces=true,                 
    numbers=left,                    
    numbersep=5pt,                  
    showspaces=false,                
    showstringspaces=false,
    showtabs=false,                  
    tabsize=2
}
\lstset{style=mystyle}

\title{\Huge \textbf{ScholarSync}\\ \vspace{1cm} \Large \textbf{Comprehensive System Architecture \&\\ Multi-Agent Workflow Documentation}}
\author{Detailed Review Preparation Notes}
\date{\today}

\begin{document}

\maketitle

\tableofcontents
\newpage

\chapter{Introduction to ScholarSync}

\section{Project Overview}
ScholarSync is an advanced, Agentic AI-powered system designed to automate the laborious process of literature reviews. Traditional AI tools (like standard ChatGPT) struggle with literature reviews because they lack persistent memory of long documents, they hallucinate facts, and they cannot map complex cross-document relationships. 

ScholarSync solves this by employing a \textbf{Multi-Agent Workflow} orchestrated by \textbf{LangGraph}. Instead of a single Large Language Model (LLM) attempting to read and summarize multiple PDFs at once, ScholarSync delegates specific tasks to specialized "agents." It combines Retrieval-Augmented Generation (RAG) with a Knowledge Graph (GraphRAG via Neo4j) to map out entities, methodologies, and findings across multiple research papers.

\section{Core Architectural Pillars}
The architecture of ScholarSync rests on five main pillars:
\begin{enumerate}
    \item \textbf{The Frontend UI:} A highly interactive, mode-switching interface that handles document uploads, visualizes agent reasoning via WebSockets, and renders a dynamic Knowledge Graph.
    \item \textbf{Document Ingestion Pipeline:} The system that parses PDFs, chunks text semantically, and generates vector embeddings.
    \item \textbf{The Multi-Agent Orchestrator (LangGraph):} The brain of the application. It manages the state machine, routing tasks between planning, extraction, validation, and synthesis agents.
    \item \textbf{GraphRAG Integration:} A Neo4j database connection that maps extracted entities (e.g., algorithms, datasets, authors) to find thematic overlaps.
    \item \textbf{Evaluation \& Quality Assurance:} A suite of post-processing scripts that score the output for hallucination risk, grounding, and source diversity.
\end{enumerate}

\newpage

\chapter{The Multi-Agent Pipeline (The Core)}

\section{What is LangGraph?}
LangGraph is a library for building stateful, multi-actor applications with LLMs. In ScholarSync, the entire research process is modeled as a cyclic graph (a state machine). The application maintains a \texttt{GraphState} dictionary that gets passed from one agent (node) to the next.

\section{\texttt{workflow/langgraph\_pipeline.py}}
This is arguably the most important file in the backend. It constructs the workflow.

\subsection{The Graph State Schema}
The state contains all the context needed for the pipeline: the user's query, the parsed metadata of the uploaded papers, subtasks generated by the manager, extractions from the workers, validation results from the checker, graph insights from Neo4j, and the final synthesized markdown report.

\subsection{Node and Edge Definitions}
The pipeline is defined with specific nodes (functions) and edges (transitions):
\begin{itemize}
    \item \textbf{Manager Node:} Decomposes the user query. Transitions to Workers.
    \item \textbf{Worker Node:} Extracts data. Transitions to GraphRAG.
    \item \textbf{GraphRAG Node:} Updates Neo4j. Transitions to Checking.
    \item \textbf{Checking Node:} Validates data. Transitions to either Synthesis (if passed) or Correction (if failed).
    \item \textbf{Correction Node:} Re-runs extraction with feedback. Transitions back to Checking (a loop).
    \item \textbf{Synthesis Node:} Writes the report. Transitions to Evaluation.
    \item \textbf{Evaluation Node:} Scores the report. Ends the pipeline.
\end{itemize}

\newpage

\chapter{Deep Dive into the Agents}

The \texttt{scholarsync/agents/} directory houses the specific logic for each agent. Each agent is essentially an LLM call wrapped with specific system prompts, structured output parsers (using Pydantic), and error handling.

\section{\texttt{manager\_agent.py}}
\textbf{Role:} The Planner.\\
\textbf{Functionality:} When a user asks a complex question (e.g., "Compare the use of CNNs vs RNNs in medical imaging"), the Manager Agent does not answer it. Instead, it analyzes the prompt and the metadata of the provided papers to create a list of \textbf{SubTasks}. For instance, Subtask 1 might be "Extract CNN methodologies from Paper A," and Subtask 2 might be "Extract RNN methodologies from Paper B."

\section{\texttt{worker\_agent.py}}
\textbf{Role:} The Extractor.\\
\textbf{Functionality:} The Worker Agent receives the subtasks. It iterates through the text chunks of the assigned papers. Using rigorous structured output formatting (JSON), it extracts:
\begin{itemize}
    \item \textbf{Entities:} Key terms, algorithms, datasets.
    \item \textbf{Relationships:} How these entities connect (e.g., "Algorithm X" [USES] "Dataset Y").
    \item \textbf{Methodology \& Findings:} Bullet points of what the paper did and what it discovered.
\end{itemize}
The workers operate in batches per paper to minimize API calls and token usage.

\section{\texttt{checking\_agent.py}}
\textbf{Role:} The Quality Gate (Anti-Hallucination mechanism).\\
\textbf{Functionality:} LLMs are prone to hallucinating facts. The Checking Agent takes the Worker Agent's extractions and compares them against the original text chunks. It assigns a validity score. If the score is below a certain threshold, it generates "Correction Prompts" explaining what the worker got wrong. The LangGraph pipeline then routes the state to the \textbf{Correction Node}, which re-runs the Worker Agent with this specific feedback.

\section{\texttt{synthesizer\_agent.py}}
\textbf{Role:} The Writer.\\
\textbf{Functionality:} Once the Checking Agent approves the extractions, the Synthesizer takes over. It receives the user's original query, the validated extractions, insights from the Knowledge Graph (cross-paper connections), and optional web search context. It writes a highly structured, academic-style markdown literature review. It ensures citations are placed correctly referencing the original PDFs.

\section{\texttt{profile\_builder.py} \& \texttt{dedup.py}}
\textbf{Role:} Post-Processing.\\
\textbf{Functionality:} \texttt{profile\_builder.py} is a deterministic script (zero LLM cost) that reformats the raw extraction data into a clean JSON structure representing a "Paper Profile" for the UI. \texttt{dedup.py} resolves duplicated information (e.g., merging similar findings).

\newpage

\chapter{Knowledge Graph \& RAG Integration}

ScholarSync doesn't just read documents sequentially; it builds a network of knowledge. This logic lives in \texttt{scholarsync/rag/} and \texttt{scholarsync/ingestion/}.

\section{Document Ingestion (\texttt{ingestion/})}
\begin{itemize}
    \item \textbf{\texttt{pdf\_loader.py}}: Uses libraries (like PyMuPDF or pdfplumber) to extract raw text from uploaded PDF files.
    \item \textbf{\texttt{chunker.py}}: LLMs have token limits. The chunker splits the massive raw text into smaller, overlapping chunks. Overlap is crucial so that context isn't lost at the boundaries of chunks.
\end{itemize}

\section{Standard RAG (\texttt{rag/embeddings.py} \& \texttt{vector\_store.py})}
For standard retrieval, text chunks are passed through an embedding model (like OpenAI's text-embedding-ada-002 or a local HuggingFace model) to create vectors. These are stored in a Vector Database. When a user asks a question, the question is embedded, and the closest matching text chunks are retrieved.

\section{GraphRAG (\texttt{rag/graph\_rag.py})}
Standard RAG fails at connecting the dots across multiple documents. GraphRAG solves this.
\begin{itemize}
    \item \textbf{Neo4j Connection:} This file establishes a connection to a Neo4j graph database.
    \item \textbf{Graph Construction:} As Worker Agents extract Entities and Relationships, \texttt{graph\_rag.py} runs Cypher queries to create nodes and edges in the database.
    \item \textbf{Session Isolation:} It uses \texttt{user\_id} scoping to ensure that one user's graph data does not pollute another user's session.
    \item \textbf{Cross-Paper Querying:} It contains functions like \texttt{query\_cross\_paper\_connections()} which ask Neo4j to find entities that appear in more than one paper. These insights are directly fed into the Synthesizer Agent to allow it to write sentences like, "Both Paper A and Paper B utilize Dataset X, but with different methodologies."
\end{itemize}

\section{\texttt{rag/entity\_normalizer.py}}
Graphs become messy if they have a node for "AI", "A.I.", and "Artificial Intelligence". This file contains logic to normalize string representations so they map to a single, unified node in Neo4j.

\newpage

\chapter{Backend Infrastructure \& the API}

The backend is built with FastAPI, providing high performance and asynchronous capabilities.

\section{\texttt{api/main.py}}
This is the application's entry point. It handles:
\begin{itemize}
    \item \textbf{CORS Configuration:} Ensuring the frontend can communicate with the backend.
    \item \textbf{Router Inclusion:} Loading the Authentication routers and Chat routers.
    \item \textbf{WebSocket Management:} The multi-agent pipeline takes time. \texttt{main.py} manages WebSockets so that as each agent completes a task, a progress message is streamed to the frontend instantly, keeping the user informed.
\end{itemize}

\section{\texttt{config/settings.py}}
A centralized configuration object (using Pydantic BaseSettings). It reads the \texttt{.env} file and provides typed access to variables like \texttt{GROQ\_API\_KEY}, \texttt{NEO4J\_URI}, and \texttt{MAX\_CORRECTION\_LOOPS}.

\section{Chat Utilities (\texttt{chat/})}
\begin{itemize}
    \item \textbf{\texttt{mode\_router.py}}: The frontend has toggles for different modes (Chat, Research, Graph). This file inspects the incoming request and routes it either to a simple LLM call or invokes the complex LangGraph pipeline.
    \item \textbf{\texttt{key\_manager.py}}: A critical stability feature. Multi-agent systems consume massive amounts of API calls quickly, often hitting provider rate limits (e.g., HTTP 429 Too Many Requests). The Key Manager holds a pool of API keys (from Groq, OpenRouter, etc.) and automatically catches rate-limit exceptions, rotating to the next available key without crashing the pipeline.
\end{itemize}

\newpage

\chapter{Evaluation, Quality Assurance, and Metrics}

ScholarSync does not just generate text; it proves its accuracy. The \texttt{scholarsync/evaluation/} directory handles post-generation scoring.

\section{\texttt{dynamic\_metrics.py}}
This script determines which metrics to calculate based on the type of query. It returns a formatted dictionary of scores that the frontend renders as a visual "Quality Report" card.

\section{\texttt{hallucination\_detector.py}}
It compares the final synthesized report against the raw source chunks retrieved during the pipeline. If the report makes claims that cannot be traced back to the source chunks, the hallucination score increases.

\section{\texttt{grounding\_checker.py}}
Similar to the hallucination detector, but it specifically checks if the report is using the designated academic citations properly. It ensures that every factual claim is "grounded" in a cited paper.

\newpage

\chapter{Authentication \& Security}

A robust, self-hosted authentication system ensures that user data and research sessions remain private.

\section{\texttt{auth/database.py} \& \texttt{models.py}}
\begin{itemize}
    \item \textbf{Database:} Uses SQLite (via SQLAlchemy) for lightweight, file-based relational storage.
    \item \textbf{Models:} Defines the \texttt{User} table (storing username, hashed password) and \texttt{Session} tables.
\end{itemize}

\section{\texttt{auth/security.py} \& \texttt{service.py}}
\begin{itemize}
    \item \textbf{Bcrypt Hashing:} Passwords are never stored in plain text. \texttt{security.py} uses \texttt{bcrypt} to hash passwords with salt before storing them.
    \item \textbf{JWT (JSON Web Tokens):} Upon successful login, the server generates a JWT. The frontend stores this token (usually in localStorage) and attaches it as a Bearer Token in the Authorization header of subsequent API requests to prove the user's identity.
\end{itemize}

\section{\texttt{auth/router.py}}
Exposes the HTTP endpoints: \texttt{/register}, \texttt{/login}, and \texttt{/me}.

\newpage

\chapter{The Frontend Interface (\texttt{chatbot-ui/})}

The UI is built using vanilla HTML, CSS, and JavaScript, optimized for speed and dynamic updates.

\section{\texttt{index.html} (Structure \& Styling)}
\begin{itemize}
    \item \textbf{Layout:} A modern sidebar layout. The sidebar contains chat history and authentication controls. The main area contains the chat bubbles and the mode selector.
    \item \textbf{CSS Variables:} Heavy use of CSS variables for theming (dark mode, accent colors), allowing for a premium, polished aesthetic with smooth transitions and glassmorphism effects.
\end{itemize}

\section{\texttt{index.html} (JavaScript Logic)}
\begin{itemize}
    \item \textbf{Mode Selection Menu:} A custom dropdown allows users to toggle Web Search, standard Chat, or Deep Research.
    \item \textbf{Streaming Logic:} Uses the Fetch API or WebSockets to read the chunked responses from the LangGraph pipeline. It dynamically creates HTML elements for "progress messages" (e.g., "Worker Agent Extracting...").
    \item \textbf{Cytoscape.js Integration:} When the pipeline completes and returns graph data, the frontend uses the Cytoscape library to render an interactive, physics-based node-and-edge visualization of the Knowledge Graph, allowing the user to visually explore how papers connect.
\end{itemize}

\newpage

\chapter{Complete Workflow Walkthrough}

To summarize the entire architecture, here is exactly what happens when a user clicks "Send" on a research query:

\begin{enumerate}
    \item \textbf{Authentication Check:} The frontend attaches the JWT. FastAPI validates it.
    \item \textbf{File Ingestion:} The user's uploaded PDFs are routed to \texttt{pdf\_loader.py}. They are parsed, chunked by \texttt{chunker.py}, and embedded.
    \item \textbf{Routing:} The frontend specifies "Deep Research". \texttt{mode\_router.py} triggers \texttt{run\_pipeline()} in \texttt{langgraph\_pipeline.py}.
    \item \textbf{Planning Phase:} The \textbf{Manager Agent} analyzes the prompt and creates Subtasks. A WebSocket message "Manager Agent: Analyzing..." is sent to the UI.
    \item \textbf{Extraction Phase:} The \textbf{Worker Agents} execute the subtasks concurrently in batches, extracting JSON data. If an API rate limit is hit, \texttt{key\_manager.py} seamlessly swaps the API key.
    \item \textbf{Graph Construction Phase:} The extracted JSON entities are sent to \texttt{graph\_rag.py}, which executes Cypher queries to populate the Neo4j database, finding overlaps between the papers.
    \item \textbf{Validation Phase:} The \textbf{Checking Agent} reviews the JSON extractions. If it spots a hallucination, it triggers the \textbf{Correction Node} to re-run the Worker for that specific chunk.
    \item \textbf{Synthesis Phase:} The \textbf{Synthesizer Agent} receives the validated extractions and the Neo4j cross-paper insights. It streams a formatted markdown report back to the frontend.
    \item \textbf{Evaluation Phase:} \texttt{dynamic\_metrics.py} calculates the hallucination and grounding scores and appends them to the final UI response.
    \item \textbf{Rendering:} The frontend renders the markdown, the metric scorecards, and draws the Cytoscape graph.
\end{enumerate}

\end{document}
